% chap09_pandas.tex - Data analysis with Pandas
% ============================================================================

\chapter{Pandas: Data Analysis}
\label{chap:pandas}

Real scientific data comes with labels, missing values, and mixed types. Pandas provides the DataFrame---a powerful table structure that handles these complexities. If NumPy is like a matrix, Pandas is like a spreadsheet with superpowers. Install pandas via:

\begin{lstlisting}
pip install pandas
\end{lstlisting}


% ----------------------------------------------------------------------------
\section{Why Pandas?}
\label{sec:why-pandas}

NumPy arrays are great for homogeneous numerical data. But experimental data often includes:

\begin{itemize}
    \item Column names (wavelength, absorbance, sample\_id)
    \item Mixed types (numbers, text, dates)
    \item Missing values
    \item Row labels (sample names, timestamps)
\end{itemize}

Pandas handles all of this with two main structures: \py{Series} (1D) and \py{DataFrame} (2D table).

\begin{lstlisting}
import pandas as pd
import numpy as np
\end{lstlisting}

% ----------------------------------------------------------------------------
\section{Series}
\label{sec:series}

A \py{Series} is a labeled one-dimensional array:

\begin{lstlisting}
# From a list
concentrations = pd.Series([0.1, 0.2, 0.3, 0.4])
print(concentrations)
# 0    0.1
# 1    0.2
# 2    0.3
# 3    0.4
# dtype: float64

# With custom labels
conc = pd.Series([0.1, 0.2, 0.3], index=["low", "medium", "high"])
print(conc["medium"])  # 0.2
\end{lstlisting}

% ----------------------------------------------------------------------------
\section{DataFrames}
\label{sec:dataframes}

A \py{DataFrame} is a labeled two-dimensional table---the workhorse of Pandas:

\begin{lstlisting}
# From a dictionary
data = {
    "sample": ["A", "B", "C", "D"],
    "concentration": [0.1, 0.2, 0.3, 0.4],
    "absorbance": [0.12, 0.25, 0.35, 0.48]
}
df = pd.DataFrame(data)
print(df)
#   sample  concentration  absorbance
# 0      A            0.1        0.12
# 1      B            0.2        0.25
# 2      C            0.3        0.35
# 3      D            0.4        0.48
\end{lstlisting}

\subsection{Reading Data from Files}

\begin{lstlisting}
# CSV (most common)
df = pd.read_csv("data.csv")

# Excel
df = pd.read_excel("data.xlsx", sheet_name="Sheet1")

# With options
df = pd.read_csv("data.csv",
                 sep=";",           # semicolon delimiter
                 skiprows=2,        # skip header rows
                 na_values=["NA", "-999"])  # missing value codes
\end{lstlisting}

\subsection{Quick Look at Data}

\begin{lstlisting}
df.head()       # first 5 rows
df.tail(3)      # last 3 rows
df.shape        # (rows, columns)
df.columns      # column names
df.dtypes       # data types
df.info()       # summary of structure
df.describe()   # statistics for numeric columns
\end{lstlisting}

% ----------------------------------------------------------------------------
\section{Selecting Data}
\label{sec:selecting}

\subsection{Selecting Columns}

\begin{lstlisting}
# Single column (returns Series)
df["concentration"]

# Multiple columns (returns DataFrame)
df[["sample", "concentration"]]

# As attribute (only for simple column names)
df.concentration
\end{lstlisting}

\subsection{Selecting Rows}

\begin{lstlisting}
# By position
df.iloc[0]      # first row
df.iloc[0:3]    # first three rows
df.iloc[[0, 2]] # rows 0 and 2

# By label
df.loc[0]       # row with label 0
df.loc[0:2]     # rows with labels 0, 1, 2 (inclusive!)

# With conditions (boolean indexing)
df[df["concentration"] > 0.2]
df[df["sample"].isin(["A", "C"])]
\end{lstlisting}

\subsection{Selecting Rows and Columns}

\begin{lstlisting}
# By position: df.iloc[rows, columns]
df.iloc[0:2, 1:3]    # first 2 rows, columns 1-2

# By label: df.loc[rows, columns]
df.loc[0:2, "concentration":"absorbance"]

# Mixed
df.loc[df["concentration"] > 0.2, ["sample", "absorbance"]]
\end{lstlisting}

\begin{tip}
Use \py{.loc} for label-based indexing (including boolean conditions) and \py{.iloc} for position-based indexing. Mixing them up causes confusion.
\end{tip}

% ----------------------------------------------------------------------------
\section{Modifying Data}
\label{sec:modifying}

\subsection{Adding Columns}

\begin{lstlisting}
# New column from calculation
df["transmittance"] = 10 ** (-df["absorbance"]) * 100

# New column from function
df["log_conc"] = np.log10(df["concentration"])
\end{lstlisting}

\subsection{Modifying Values}

\begin{lstlisting}
# Set specific cells
df.loc[0, "concentration"] = 0.15

# Set based on condition
df.loc[df["absorbance"] > 0.4, "quality"] = "high"

# Apply function to column
df["sample"] = df["sample"].str.lower()
\end{lstlisting}

\subsection{Renaming}

\begin{lstlisting}
df.rename(columns={"concentration": "conc_M"}, inplace=True)
\end{lstlisting}

\subsection{Dropping}

Apply the following then inspect your dataframe.

\begin{lstlisting}
# Drop columns
df.drop(columns=["transmittance"], inplace=True)

# Drop rows
df.drop(index=[0, 1], inplace=True)
\end{lstlisting}

What has happed to the data?

% ----------------------------------------------------------------------------
\section{Handling Missing Data}
\label{sec:missing}

Missing values appear as \py{NaN} (Not a Number):

\begin{lstlisting}
# Check for missing values
df.isna()           # boolean mask
df.isna().sum()     # count per column

# Drop rows with any missing values
df.dropna()

# Drop rows where specific columns are missing
df.dropna(subset=["absorbance"])

# Fill missing values
df.fillna(0)                    # fill with 0
df.fillna(df.mean())            # fill with column means
df.fillna(method="ffill")       # forward fill
\end{lstlisting}

\begin{labexample}
Handling missing readings in experimental data:

\begin{lstlisting}
# Data with missing values (-999 indicates error)
data = pd.read_csv("measurements.csv")
data.replace(-999, np.nan, inplace=True)

# Report missing data
missing = data.isna().sum()
print(f"Missing values:\n{missing}")

# Remove rows with missing absorbance
clean = data.dropna(subset=["absorbance"])
print(f"Kept {len(clean)} of {len(data)} rows")
\end{lstlisting}
\end{labexample}

% ----------------------------------------------------------------------------
\section{Aggregation and Grouping}
\label{sec:aggregation}

\subsection{Basic Statistics}

\begin{lstlisting}
df["absorbance"].mean()
df["absorbance"].std()
df["absorbance"].min()
df["absorbance"].max()
df["concentration"].sum()
df["sample"].value_counts()  # count unique values
\end{lstlisting}

\subsection{Group By}

Split data into groups and compute statistics:

\begin{lstlisting}
# Sample data
experiments = pd.DataFrame({
    "treatment": ["A", "A", "B", "B", "A", "B"],
    "replicate": [1, 2, 1, 2, 3, 3],
    "yield": [23, 25, 31, 33, 24, 35]
})

# Group and aggregate
grouped = experiments.groupby("treatment")
print(grouped["yield"].mean())
# treatment
# A    24.0
# B    33.0

# Multiple aggregations
experiments.groupby("treatment")["yield"].agg(["mean", "std", "count"])
\end{lstlisting}

\begin{labexample}
Analyzing replicate measurements:

\begin{lstlisting}
data = pd.DataFrame({
    "sample": ["S1", "S1", "S1", "S2", "S2", "S2"],
    "absorbance": [0.45, 0.47, 0.44, 0.82, 0.85, 0.81]
})

summary = data.groupby("sample")["absorbance"].agg(
    mean="mean",
    std="std",
    n="count"
)
summary["sem"] = summary["std"] / np.sqrt(summary["n"])
print(summary)
\end{lstlisting}
\end{labexample}

% ----------------------------------------------------------------------------
\section{Sorting}
\label{sec:sorting}

\begin{lstlisting}
# Sort by column
df.sort_values("absorbance")
df.sort_values("absorbance", ascending=False)

# Sort by multiple columns
df.sort_values(["treatment", "concentration"])

# Sort index
df.sort_index()
\end{lstlisting}

% ----------------------------------------------------------------------------
\section{Merging DataFrames}
\label{sec:merging}

Combine data from multiple sources:

\begin{lstlisting}
# Sample data
samples = pd.DataFrame({
    "sample_id": ["A", "B", "C"],
    "concentration": [0.1, 0.2, 0.3]
})

metadata = pd.DataFrame({
    "sample_id": ["A", "B", "D"],
    "temperature": [25, 30, 25]
})

# Merge on common column
merged = pd.merge(samples, metadata, on="sample_id", how="inner")
# Only keeps rows present in both (A, B)

# Different join types
pd.merge(samples, metadata, how="left")   # keep all from samples
pd.merge(samples, metadata, how="right")  # keep all from metadata
pd.merge(samples, metadata, how="outer")  # keep all from both
\end{lstlisting}

% ----------------------------------------------------------------------------
\section{Pivot Tables}
\label{sec:pivot}

Reshape data for analysis:

\begin{lstlisting}
data = pd.DataFrame({
    "date": ["Mon", "Mon", "Tue", "Tue"],
    "sample": ["A", "B", "A", "B"],
    "value": [1.2, 2.3, 1.4, 2.5]
})

# Pivot: rows=dates, columns=samples, values=value
pivot = data.pivot(index="date", columns="sample", values="value")
#        A    B
# date
# Mon  1.2  2.3
# Tue  1.4  2.5
\end{lstlisting}

% ----------------------------------------------------------------------------
\section{Applying Functions}
\label{sec:applying}

\begin{lstlisting}
# Apply to entire column
df["absorbance"].apply(lambda x: x * 1000)  # convert to mAU

# Apply to each row
def classify_sample(row):
    if row["absorbance"] > 0.5:
        return "high"
    return "low"

df["category"] = df.apply(classify_sample, axis=1)

# Vectorized operations (faster)
df["category"] = np.where(df["absorbance"] > 0.5, "high", "low")
\end{lstlisting}

% ----------------------------------------------------------------------------
\section{Saving Data}
\label{sec:saving}

\begin{lstlisting}
# To CSV
df.to_csv("results.csv", index=False)

# To Excel
df.to_excel("results.xlsx", sheet_name="Data", index=False)

# Multiple sheets
with pd.ExcelWriter("report.xlsx") as writer:
    raw_data.to_excel(writer, sheet_name="Raw")
    summary.to_excel(writer, sheet_name="Summary")
\end{lstlisting}

% ----------------------------------------------------------------------------
\section{Working with Dates and Times}
\label{sec:datetime}

Scientific data is often timestamped: instrument readings, reaction monitoring, environmental measurements. Python and Pandas provide powerful tools for handling temporal data.

\subsection{Python's datetime Module}

The standard library \py{datetime} module handles dates and times:

\begin{lstlisting}
from datetime import datetime, date, time, timedelta

# Current date and time
now = datetime.now()
print(now)  # 2024-03-15 14:30:45.123456

# Create specific dates/times
experiment_date = date(2024, 3, 15)
start_time = time(9, 30, 0)  # 9:30:00 AM
measurement = datetime(2024, 3, 15, 14, 30, 45)

# Access components
print(measurement.year)    # 2024
print(measurement.month)   # 3
print(measurement.hour)    # 14
print(measurement.weekday())  # 4 (Friday, where Monday=0)
\end{lstlisting}

\subsubsection{Formatting and Parsing}

Convert between datetime objects and strings using format codes:

\begin{lstlisting}
from datetime import datetime

# Format datetime as string
dt = datetime(2024, 3, 15, 14, 30)
formatted = dt.strftime("%Y-%m-%d %H:%M")
print(formatted)  # "2024-03-15 14:30"

# Parse string to datetime
text = "15/03/2024 09:45"
parsed = datetime.strptime(text, "%d/%m/%Y %H:%M")
print(parsed)  # 2024-03-15 09:45:00
\end{lstlisting}

Common format codes: \py{\%Y} (4-digit year), \py{\%m} (month), \py{\%d} (day), \py{\%H} (24-hour), \py{\%M} (minute), \py{\%S} (second).

\subsubsection{Time Differences}

Use \py{timedelta} for durations and arithmetic:

\begin{lstlisting}
from datetime import datetime, timedelta

start = datetime(2024, 3, 15, 9, 0)
end = datetime(2024, 3, 15, 17, 30)

duration = end - start
print(duration)              # 8:30:00
print(duration.total_seconds())  # 30600.0

# Add time to a datetime
next_reading = start + timedelta(hours=2, minutes=30)
print(next_reading)  # 2024-03-15 11:30:00
\end{lstlisting}

\subsection{DateTime in Pandas}

Pandas extends datetime handling for working with columns of timestamps.

\subsubsection{Converting to DateTime}

Use \py{pd.to\_datetime()} to parse date columns:

\begin{lstlisting}
import pandas as pd

# From various string formats
df = pd.DataFrame({
    "date_str": ["2024-03-15", "2024-03-16", "2024-03-17"],
    "value": [1.2, 1.5, 1.3]
})
df["date"] = pd.to_datetime(df["date_str"])

# Pandas often infers format automatically
dates = pd.to_datetime(["March 15, 2024", "March 16, 2024"])

# Specify format for speed or unusual formats
df["date"] = pd.to_datetime(df["date_str"], format="%Y-%m-%d")
\end{lstlisting}

\subsubsection{The .dt Accessor}

Access datetime components using \py{.dt}:

\begin{lstlisting}
df["date"] = pd.to_datetime(df["date_str"])

df["year"] = df["date"].dt.year
df["month"] = df["date"].dt.month
df["day_of_week"] = df["date"].dt.dayofweek  # Monday=0
df["hour"] = df["date"].dt.hour

# Useful for filtering
morning_data = df[df["date"].dt.hour < 12]
march_data = df[df["date"].dt.month == 3]
\end{lstlisting}

\subsubsection{DateTime as Index}

Setting timestamps as the index enables powerful time series operations:

\begin{lstlisting}
# Create time series data
data = pd.DataFrame({
    "timestamp": pd.date_range("2024-03-15", periods=24, freq="h"),
    "temperature": [20.1, 20.3, 20.0, 19.8, 19.5, 19.2,
                   19.0, 19.5, 21.0, 23.5, 25.2, 26.8,
                   27.5, 28.0, 27.8, 27.0, 25.5, 24.0,
                   22.5, 21.5, 21.0, 20.5, 20.2, 20.0]
})
data = data.set_index("timestamp")

# Select by date/time
march_15_afternoon = data.loc["2024-03-15 12:00":"2024-03-15 18:00"]

# Resample: convert hourly to 6-hour averages
six_hourly = data.resample("6h").mean()
print(six_hourly)
\end{lstlisting}

\subsubsection{Resampling Time Series}

Resampling converts data between time frequencies:

\begin{lstlisting}
# Common frequency codes:
# "D" = daily, "h" = hourly, "min" = minute
# "W" = weekly, "ME" = month end, "YE" = year end

# Downsample: reduce frequency (aggregate)
daily_mean = hourly_data.resample("D").mean()
daily_max = hourly_data.resample("D").max()

# Multiple aggregations
daily_stats = hourly_data.resample("D").agg(["mean", "min", "max"])
\end{lstlisting}

\begin{labexample}
Analyzing timestamped instrument readings:

\begin{lstlisting}
import pandas as pd

# Load data with timestamps
data = pd.read_csv("instrument_log.csv")
data["timestamp"] = pd.to_datetime(data["timestamp"])
data = data.set_index("timestamp")

# Find readings from specific time window
morning = data.between_time("08:00", "12:00")

# Calculate hourly statistics
hourly = data["reading"].resample("h").agg(["mean", "std", "count"])
print(hourly.head())

# Find days with high variance (potential issues)
daily_std = data["reading"].resample("D").std()
problem_days = daily_std[daily_std > daily_std.mean() + 2*daily_std.std()]
print(f"High-variance days: {len(problem_days)}")
\end{lstlisting}
\end{labexample}

% ----------------------------------------------------------------------------
\section{A Complete Analysis Example}
\label{sec:pandas-example}

\begin{lstlisting}
import pandas as pd
import numpy as np

# Load calibration data
calib = pd.read_csv("calibration.csv")
print("Calibration data:")
print(calib.head())

# Check for missing values
if calib.isna().any().any():
    print("Warning: missing values found")
    calib = calib.dropna()

# Calculate statistics by standard
stats = calib.groupby("standard_id")["absorbance"].agg(
    ["mean", "std", "count"]
)

# Load unknown samples
unknowns = pd.read_csv("unknowns.csv")

# Calculate concentration from calibration curve
# (assuming linear: A = m*C + b, pre-calculated)
slope = 1.5
intercept = 0.02
unknowns["concentration"] = (unknowns["absorbance"] - intercept) / slope

# Add quality flags
unknowns["in_range"] = (
    (unknowns["concentration"] >= calib["concentration"].min()) &
    (unknowns["concentration"] <= calib["concentration"].max())
)

# Save results
unknowns.to_csv("results.csv", index=False)
print(f"\nProcessed {len(unknowns)} samples")
print(f"Samples in calibration range: {unknowns['in_range'].sum()}")
\end{lstlisting}

\begin{tip}
As for numpy, you are highly encourage to find and peruse the documentation for pandas in order to find out more about it's features. For example, it offers plotting funcitonality, which can be useful but is not essential to cover as we will cover plotting in detail next.
\end{tip}

% ----------------------------------------------------------------------------
\section{Chapter Summary}
\label{sec:pandas-summary}

Pandas makes data analysis manageable:

\begin{itemize}
    \item \textbf{DataFrame:} Labeled 2D table with mixed types
    \item \textbf{Reading:} \py{pd.read\_csv()}, \py{pd.read\_excel()}
    \item \textbf{Selection:} \py{df["col"]}, \py{df.loc[]}, \py{df.iloc[]}
    \item \textbf{Missing data:} \py{df.isna()}, \py{df.dropna()}, \py{df.fillna()}
    \item \textbf{Grouping:} \py{df.groupby("col").agg()}
    \item \textbf{Merging:} \py{pd.merge(df1, df2, on="key")}
    \item \textbf{Saving:} \py{df.to\_csv()}, \py{df.to\_excel()}
\end{itemize}

\begin{ponder}
\begin{enumerate}
    \item What's the difference between \py{df.loc[0]} and \py{df.iloc[0]}?
    \item How would you find rows where column A is greater than 5 AND column B equals ``yes''?
    \item When would you use \py{groupby()} vs. just calculating \py{.mean()} on a column?
    \item How do the different merge types (inner, left, right, outer) differ?
\end{enumerate}
\end{ponder}
